\chapter{团队与组织——谁在训练这些模型？}
\label{ch:org}
%% ============================================================

训练一个 frontier LLM 不仅是技术挑战，
更是一个复杂的组织工程。
它需要多个专业团队的紧密协作，
每个团队都面临独特的技术和管理挑战。
一个常见的误解是认为 LLM 训练的核心是算法研究——
事实上，工程执行占据了整个工作量的 80\% 以上。

\section{核心团队}

\subsection{数据团队：被低估的英雄}

数据团队是 LLM 训练中最被低估、但影响力最大的团队。
一个在业界广泛流传的经验法则是：
数据工程占据了整个训练工作量的 60--80\%。
模型架构可以从论文中复制，训练框架可以用开源工具，
但高质量的数据流水线必须从零构建。

数据团队的核心角色包括：

\begin{itemize}[noitemsep]
  \item \textbf{数据管道工程师}（Data Pipeline Engineers）：
        构建和维护从网页爬取到最终训练数据的全流程基础设施。
        这包括爬虫集群管理、数据清洗 pipeline（通常基于 Spark 或 Dask 构建）、
        去重系统（MinHash + LSH 在 PB 级数据上的高效实现）、
        以及数据版本管理（类似 Git LFS 但针对 TB 级数据集的版本控制）。
        这些工程师需要同时精通分布式系统和 NLP。
  \item \textbf{数据质量分析师}（Data Quality Analysts）：
        定义和维护数据过滤标准——什么样的文本是「高质量」的？
        这不是一个纯技术问题。
        perplexity 过低可能意味着模板化文本，过高可能意味着乱码。
        重复率、特殊字符比例、段落结构完整性——
        每一个过滤规则都需要在大量样本上验证，
        确保不会误杀有价值的数据。
  \item \textbf{语言学家}（Linguists）：
        对于多语言模型，语言学家的角色至关重要。
        分词器的设计需要考虑不同语言的形态学特征（中文的字/词粒度、
        日语的假名与汉字混合、阿拉伯语的从右到左书写和词根系统）。
        语言学家还负责评估不同语言的数据质量标准——
        适用于英语的 perplexity 过滤阈值不能直接套用到中文。
  \item \textbf{标注管理者}（Annotation Managers）：
        协调人类标注员团队，管理标注质量。
        在 SFT 和 RLHF 阶段，高质量的人类标注数据是稀缺资源。
        标注管理者需要设计标注指南、培训标注员、
        建立质量控制流程（inter-annotator agreement 监控、
        抽样审核、异常检测），
        并在成本和质量之间寻找平衡。
\end{itemize}

在一个主要 AI 实验室中，数据团队的规模通常在 10--30 人。
DeepSeek 在其技术报告~\cite{deepseekv3}中披露了
对数据配比的大量实验——网页、代码、学术论文、数学文本的最优比例
是通过在小模型上进行数百次 proxy 实验确定的。
这种工作量是纯算法研究无法覆盖的。

\begin{whybox}[为什么「数据质量 > 数据数量」？]
FineWeb~\cite{penedo2024fineweb} 的实验表明，
经过精细过滤的 1.3T tokens 数据集训练出的模型，
在多个基准测试上优于使用 4T tokens 低质量数据训练的模型。
直觉解释是：低质量数据不仅「浪费」了训练步数，
还可能引入噪声模式（如 HTML 标签、广告文本、机器翻译痕迹），
这些模式在推理时会以微妙的方式降低输出质量。
这就是为什么数据团队的工作如此关键。
\end{whybox}

\subsection{基础设施团队：让数千块 GPU 协同工作}

基础设施（Infra）团队负责 GPU 集群的管理与优化。
这个团队面临的核心挑战可以用一个数字概括：
在一个 16,384 块 GPU 的集群中，
假设每块 GPU 每天的故障率为 0.1\%，
那么平均每天会发生 $16384 \times 0.001 \approx 16$ 次 GPU 故障。
这意味着\textbf{平均每 90 分钟就有一块 GPU 出问题}。

Infra 团队的核心职责包括：

\begin{itemize}[noitemsep]
  \item \textbf{集群管理与任务调度}：
        使用 SLURM（传统 HPC 集群的标配）或 Kubernetes
        管理训练任务的调度。
        关键挑战是资源隔离——当多个训练任务共享集群时，
        如何确保网络带宽不互相干扰？
        如何在节点故障时自动重新调度而不中断其他任务？
  \item \textbf{故障检测与自动恢复}：
        这是 Infra 团队最核心的能力。
        自动检测 GPU 故障（CUDA error、ECC 错误、NVLink 断开）、
        网络故障（InfiniBand/RoCE 链路中断）、
        存储故障（NFS/Lustre 挂载点失效）。
        检测到故障后，系统需要自动隔离故障节点、
        将训练任务迁移到备用节点、
        从最近的 checkpoint 恢复训练——整个过程最好在 5--10 分钟内完成。
  \item \textbf{Checkpoint 管理}：
        对于 70B 参数的模型，每个完整训练 checkpoint 的大小约为
        $70 \times 10^9 \times 8 = 560\,\mathrm{GB}$
        （包含 BF16 参数、FP32 master weights 和 Adam 优化器状态 $m$/$v$，
        详见附录）。
        如果每 1000 步保存一次，保留最近 10 个 checkpoint，
        就需要 5.6\,TB 的高速存储空间。
        Infra 团队需要设计 checkpoint 的写入策略
        （异步写入，不阻塞训练）和存储架构
        （通常使用 Lustre 或 GPFS 这样的并行文件系统，
        提供数十 GB/s 的聚合带宽）。
  \item \textbf{网络监控与优化}：
        分布式训练的通信量巨大。
        在 tensor parallelism 中，
        每个 forward 和 backward pass 都需要 all-reduce 操作。
        Infra 团队需要监控网络拓扑中的热点、
        检测慢节点（stragglers）、
        优化 NCCL 的通信参数（如环形 vs.\ 树形 all-reduce、
        通信与计算的重叠度）。
\end{itemize}

\begin{corebox}[DeepSeek 的基础设施智慧]
DeepSeek~\cite{deepseekv3} 在 H800 集群（受出口管制限制，
NVLink 带宽低于 H100）上训练 V3 模型，
通过软件层面的创新弥补了硬件劣势：
DualPipe 通过更精细的计算-通信重叠减少了 pipeline bubble，
DeepEP 则为 MoE 的 expert parallelism
设计了专用的 all-to-all 通信内核。
他们实现了 ``zero irrecoverable loss spikes''，
即整个训练过程中没有一次不可恢复的 loss 异常——
这需要极其成熟的容错基础设施。
\end{corebox}

\subsection{预训练研究团队：训练的「配方师」}

预训练研究团队负责回答一系列关键问题：
模型应该多大（多少层？多少注意力头？多大的 hidden dimension？）；
训练应该怎么做（学习率多少？batch size 多大？warmup 多长？数据怎么混合？）。

这个团队的工作方法是\textbf{先在小规模上验证，再推广到大规模}。
典型的工作流程：

\begin{enumerate}[noitemsep]
  \item 在 1B--7B 的 proxy model 上做超参数搜索，
        每组实验训练数小时到数天。
  \item 利用 Scaling Laws~\cite{kaplan2020scaling,hoffmann2022chinchilla}
        将小规模的最优配方推算到目标规模（如 70B 或 671B）。
  \item 在目标规模上做少量验证实验（训练数千步，
        观察 loss 曲线是否符合预期）。
  \item 确认配方后启动完整训练，通常持续数周到数月。
\end{enumerate}

「训练配方」（training recipe）是每个实验室的核心知识资产。
它包括但不限于：
\textbf{学习率调度}（cosine decay 的起点、终点和 warmup 步数）、
\textbf{batch size 调度}（从小 batch 逐渐增大，
还是全程使用固定 batch size？）、
\textbf{数据配比调度}（训练前期多用网页数据建立通用能力，
后期增加代码和数学数据提升推理能力）、
\textbf{长上下文扩展策略}（何时从 4K 切换到 128K，
RoPE 频率基数如何调整）。

LLaMA 3~\cite{dubey2024llama3} 的训练报告
详细披露了他们的训练配方：
学习率峰值 $1.5 \times 10^{-4}$，
cosine 衰减到峰值的 $1/30$，
warmup 8000 步，batch size 从 4M tokens 逐步增加到 16M tokens。
这些数字看似简单，但每一个都是大量实验的结果。

\subsection{后训练/对齐团队：定义「好」的含义}

后训练团队的工作是将一个「博学但不可控」的 base model
转变为一个「有用、无害、诚实」的 assistant。
这是整个训练流程中最依赖人类判断的环节。

核心职责包括：

\begin{itemize}[noitemsep]
  \item \textbf{RLHF/DPO 训练管道}：
        设计和实现偏好优化的训练流程。
        收集人类偏好数据（A 回答 vs.\ B 回答，哪个更好？），
        训练 reward model 或直接使用 DPO/GRPO 优化。
        关键挑战是偏好数据的一致性——
        不同标注员对「更好」的定义可能不一致。
  \item \textbf{Constitutional AI 设计}：
        受 Anthropic~\cite{bai2022constitutional} 启发，
        一些团队使用 AI 自身来生成和评估训练数据，
        基于一组预定义的「宪法原则」。
        这减少了对人类标注的依赖，
        但需要精心设计原则体系——
        什么是「有害的」？什么程度的帮助是「过度的」？
  \item \textbf{Red teaming 与安全评估}：
        组建专业的红队，系统性地测试模型的安全边界。
        红队会尝试越狱攻击（jailbreaking）、
        提示注入（prompt injection）、
        社会工程学手段（伪装成合法用户请求危险信息）。
        每一个发现的漏洞都需要通过额外的对齐训练来修复，
        然后再次测试——这是一个持续的迭代过程。
  \item \textbf{推理能力训练}：
        对于 reasoning model（如 DeepSeek-R1~\cite{deepseekr1}），
        后训练团队需要设计 RL 训练流程，
        让模型学会长链推理（chain-of-thought）、
        自我验证和错误修正。
        这需要精心设计的 reward signal——
        数学问题可以用答案正确性作为 reward，
        但开放性问题的 reward 设计要困难得多。
\end{itemize}

\begin{warnbox}[对齐的核心困难]
后训练团队面临的最根本挑战是：
\textbf{「有用、无害、诚实」这三个目标之间存在张力。}
过度追求无害会导致模型过度拒绝合理请求（over-refusal）；
过度追求有用会导致模型协助有害行为；
过度追求诚实会导致模型在不确定时拒绝给出有用的近似回答。
找到这三者的平衡点，本质上是一个价值观设计问题，
没有唯一正确的答案。
\end{warnbox}

\subsection{评估团队：度量你所关心的}

评估团队的工作是回答一个看似简单但极其困难的问题：
\textbf{模型到底有多好？}

核心职责包括：

\begin{itemize}[noitemsep]
  \item \textbf{Benchmark 套件管理}：
        维护覆盖多个能力维度的评估基准集合——
        数学推理（MATH、GSM8K）、代码生成（HumanEval、SWE-bench）、
        知识问答（MMLU）、指令遵循（IFEval）、
        长文本理解、多语言能力等。
        每个 benchmark 都需要持续维护——
        防止数据泄露（训练数据中包含了测试集）、
        更新过时的题目、处理评估格式变化。
  \item \textbf{内部评估管道}：
        构建自动化的评估基础设施。
        每次生成新的 checkpoint，
        系统需要自动在数十个 benchmark 上运行评估，
        生成对比报告。这本身就需要大量 GPU 资源——
        对 70B 模型在 MMLU 的 14K+ 题目上做推理
        可能需要数小时。
        高效的评估管道（如 lm-evaluation-harness 或 Lighteval）
        和评估集群的合理调度至关重要。
  \item \textbf{人类评估与 A/B 测试}：
        Benchmark 分数不能完全反映用户体验。
        评估团队还需要组织人类评估——
        让真实用户或专业评估者对比不同版本的模型输出，
        给出偏好判断。
        Chatbot Arena（LMSYS）的 Elo 评分系统
        是目前最受认可的人类评估方法。
  \item \textbf{评估的局限性}：
        评估团队最重要的认知是：
        \textbf{没有完美的评估}。
        模型可能在所有现有 benchmark 上表现优异，
        但在真实场景中出现令人意外的失败。
        评估团队需要持续设计新的评估方法，
        特别是针对 benchmark 无法捕捉的能力
        （如长期规划、常识推理、文化理解）。
\end{itemize}

\subsection{推理/部署团队：连接模型与用户}

推理/部署团队负责将训练好的模型变成用户可以使用的产品。
这个团队的工作直接决定了用户体验和运营成本。

核心职责包括：

\begin{itemize}[noitemsep]
  \item \textbf{服务架构设计}：
        API 网关 → 负载均衡器 → 推理集群 → 模型注册中心。
        需要处理突发流量（某个功能在社交媒体上火了，
        请求量可能在数分钟内增长 10 倍）、
        优雅降级（部分节点故障时不影响整体可用性）、
        自动扩缩容（根据负载动态调整推理服务器数量）。
  \item \textbf{延迟与成本优化}：
        用户期望的首 token 延迟（TTFT）通常在 500ms 以内，
        生成速度期望在 30--100 tokens/s。
        同时，推理成本直接影响定价和利润。
        部署团队需要在量化级别（FP8 vs.\ INT4）、
        批处理策略（continuous batching 的 max batch size）、
        缓存策略（prefix caching 的命中率优化）
        之间找到最优平衡。
  \item \textbf{SLA 管理}：
        商业 API 需要保证 99.9\% 以上的可用性
        （每月最多 43 分钟停机）。
        这需要多区域部署、灰度发布（新版本先在小流量上验证）、
        自动回滚机制。
  \item \textbf{滥用检测}：
        监控异常使用模式——大量请求尝试绕过安全限制、
        token farming（批量生成内容）、API key 泄露等。
        需要实时的异常检测系统和响应流程。
\end{itemize}

\section{组织模式：从实验室到公司}

观察最成功的 LLM 团队，
一个共同特点是\textbf{研究与工程的深度融合}——
写论文的人和写训练代码的人往往是同一批人。
但不同组织在规模、文化和战略上有显著差异。

\subsection{主要玩家深度解析}
\label{sec:org-labs}

\textbf{OpenAI}（2026 年初已超过 7,000 人）
经历了 AI 领域最戏剧化的组织演变。
从 2015 年的非营利研究实验室，到 2019 年的 ``capped profit'' 混合结构，
再到 2025 年底正式转型为营利性公司——
每一步都伴随着巨大的争议。
微软累计投资超过 130 亿美元，
而 2025 年 OpenAI 在日本软银领投下完成了 400 亿美元的融资，
估值达到 3,000 亿美元。
组织上，OpenAI 在 2025--2026 年经历了显著的人才流动：
联合创始人 Ilya Sutskever 离开创立了 Safe Superintelligence Inc.（SSI），
多位安全团队核心成员出走。
CEO Sam Altman 将 2025 年 8 月任命前 Instagram 负责人
Fidji Simo 为 COO 并在 2026 年 2 月升任 CEO（应用）,
自己则转为专注 AGI 研究和战略。
OpenAI 的组织高度垂直整合——
从基础研究（GPT 系列、o 系列推理模型）到消费产品（ChatGPT，月活超 4 亿）、
到企业 API、到硬件（自研芯片计划与 Broadcom 合作），
全部在一个屋檐下。
内部团队结构包括：
\textbf{Foundation Model} 团队（预训练与 scaling）、
\textbf{Post-Training} 团队（RLHF、推理训练）、
\textbf{Safety Systems} 团队（红队与对齐）、
\textbf{Applied AI/ChatGPT} 团队（产品工程）、
以及新成立的 \textbf{Deep Research} 团队（长链推理与信息综合）。

\textbf{Anthropic}（2025 年约 1,500 人，2026 年持续快速增长）
由前 OpenAI 研究副总裁 Dario Amodei 和他的姐姐 Daniela Amodei 于 2021 年创立，
核心使命是 AI 安全研究。
Anthropic 的技术路线以 Constitutional AI~\cite{bai2022constitutional} 为基石——
用 AI 自身而非纯人工来完成对齐工作。
在 Claude 模型系列中，Anthropic 展现了独特的产品策略：
2025 年推出了 Claude 3.7 Sonnet（首个混合推理模型，
同时支持快速回答和 extended thinking 深度推理），
2026 年初发布了 Claude Opus 4 和 Claude Sonnet 4（代号 Kryptonite 和 Borealis），
在编程和推理任务上取得了领先地位。
组织特点是安全团队在决策中有异常大的话语权——
新能力的部署需要通过 Responsible Scaling Policy（RSP）评估，
模型在发布前需要通过一系列 AI Safety Levels（ASL）评估。
Anthropic 在 2025 年完成了 Google、Salesforce 和
Amazon 领投的多轮融资，估值超过 600 亿美元。
值得注意的产品创新包括：
Claude Code（CLI 级别的自主编程 Agent）、
Computer Use（通过截图和键鼠操作桌面应用）、
以及 Model Context Protocol（MCP，Agent 工具交互的开放标准）。

\textbf{Google DeepMind}（数千人，是规模最大的 AI 研究组织）
2023 年由 Google Brain 和 DeepMind 合并而成，
由 Demis Hassabis 领导。
拥有最多的计算资源（数十万块 TPU 组成的超级集群）
和最大的研究团队，
但组织规模也带来了协调成本。
Gemini 系列是其旗舰产品：
Gemini 2.0 Flash（2024 年底）引入了原生工具使用和代码执行，
Gemini 2.5 Pro（2025 年初）成为首个在 LMSYS Chatbot Arena 登顶的 Google 模型，
而 Gemini 3（预计 2026 年中发布）
将进一步推进原生多模态和长上下文能力。
Google DeepMind 的独特优势是
\textbf{端到端的基础设施控制}——
从芯片设计（TPU）到网络架构到训练框架（JAX/Flax）
到云服务（Google Cloud）再到消费产品
（Gemini App、Google Search 的 AI Overviews），
全栈自研。
但合并后的组织需要持续在「学术自由」（DeepMind 的传统）
和「产品交付」（Google 的商业需求）之间寻找平衡。
多个产品团队（Gemini App、Workspace AI、Cloud AI）
有时会出现内部竞争。

\textbf{Meta FAIR 及 GenAI 团队}（合计 1,000+ 人）
Mark Zuckerberg 做出了 AI 行业最大胆的战略选择：
将 LLaMA~\cite{touvron2023llama,dubey2024llama3} 系列开源。
这不是慈善——开源是一种竞争策略：
当你的模型成为行业标准，围绕它的生态系统（工具、微调模型、应用）
就成为你的护城河。Meta 不靠卖 API 赚钱，
而是通过 AI 能力提升其广告和社交产品的收入。
LLaMA 4 系列（2025 年 4 月发布，包括 Scout 109B/17B-active
和 Maverick 400B/17B-active）采用了 MoE 架构，
是 Meta 开源的首个 MoE 模型。
组织上，Meta AI 分为两个核心部门：
FAIR（Fundamental AI Research，基础研究）
和 GenAI（生成式 AI 产品工程）。
Yann LeCun 仍担任 Chief AI Scientist，
负责 FAIR 的长期研究方向。

\textbf{DeepSeek}（约 150--200 名核心研究员和工程师）
背靠量化对冲基金幻方量化（High-Flyer Capital），
DeepSeek 是「小团队大影响」的典范。
他们用不到 Meta FAIR 十分之一的人员，
在 DeepSeek-V3~\cite{deepseekv3} 和 R1~\cite{deepseekr1}
上达到了 GPT-4 级别的性能。
关键因素：极致的工程效率（MLA、DeepEP 等创新）、
扁平的组织结构（减少管理层级，加速决策）、
以及量化基金背景带来的对效率的极度重视。
DeepSeek 的团队结构高度集成——
没有独立的「数据团队」「Infra 团队」「研究团队」的明确边界，
研究员同时负责算法设计和底层系统优化。
创始人梁文锋（Wenfeng Liang）的管理风格
强调技术深度和第一性原理思考，
团队成员大多来自顶尖高校和量化行业。
2025 年 1 月 DeepSeek-R1 的发布
（被称为 ``China's Sputnik moment''）
引发了全球对中国 AI 能力的重新评估。

\textbf{Mistral}（约 600--700 人，2025 年）
欧洲最重要的 AI 公司，创始团队来自 Google DeepMind 和 Meta。
Mixtral~\cite{jiang2024mixtral} 证明了
一个精干的团队可以在 MoE 架构上做出出色的工作。
Mistral 的策略是专注于效率——
不追求最大的模型，而是在给定规模下做到最好。
2025 年，Mistral 推出了 Mistral Large 2 和 Pixtral（多模态），
同时为欧洲企业提供了符合 EU AI Act 合规要求的定制部署方案，
这成为其在欧洲市场的重要竞争优势。

\textbf{其他值得关注的组织}：
\textbf{xAI}（Elon Musk 创立，2024 年部署了超过 10 万块 H100 组成的
``Colossus'' 集群，推出了 Grok 3 系列模型）；
\textbf{阿里巴巴通义}（Qwen 系列，特别是 Qwen 3
在开源多语言模型中表现突出~\cite{qwen2025qwen3}）；
\textbf{AI2}（Allen Institute for AI，OLMo 系列是
真正意义上的完全开源模型——数据、代码、权重全部公开）；
\textbf{Cohere}（专注企业市场，Command R+ 系列针对 RAG 和企业搜索优化）。

\begin{keybox}[「小团队，大影响」的启示]
DeepSeek（150--200 名研究员）的成功揭示了一个重要规律：
在 LLM 开发中，
\textbf{团队的工程文化和技术深度
比纯粹的人员规模更重要}。
小团队的优势包括：
（1）更快的决策速度——没有多层审批，
实验想法从提出到验证可以在数天内完成；
（2）更强的 ownership——每个人负责更大的模块，
减少了沟通开销和责任模糊；
（3）更高的人才密度——小团队可以只招最优秀的人，
大团队不可避免地需要大量普通工程师做基础工作。
Mistral 在 2023 年创立初期也以 50--80 人的精干团队
产出了 Mixtral 等有影响力的模型，
印证了同样的规律。
但随着商业化需求增长，Mistral 已扩展至 600--700 人（2025 年），
说明规模化终究是从研究到产品的必经之路。
一个有趣的对比：OpenAI 的 7,000+ 人产出了 GPT-4o/o3，
DeepSeek 的 200 人产出了 V3/R1——
在核心模型质量上，两者差距远小于团队规模的差距。
\end{keybox}

\section{训练成本的演变：2024--2026}
\label{sec:training-cost}

训练 frontier LLM 的成本正在快速上升，
但效率的提升也同样惊人。
理解这个双重趋势对于评估 AI 行业的可持续性至关重要。

\subsection{成本的量级}

截至 2026 年初，已知的训练成本估算如下：

\begin{table}[H]
\centering\small
\caption{Frontier 模型训练成本估算（2024--2026）}
\begin{tabularx}{\linewidth}{l l l X}
\toprule
\textbf{模型} & \textbf{训练时间} & \textbf{估算成本} & \textbf{备注} \\
\midrule
GPT-4 (2023) & $\sim$3 个月 & \$78M--100M+ &
  早期估算；含 A100 集群的 GPU 租赁成本 \\
Gemini Ultra (2023) & 数月 & \$191M（估算） &
  使用 Google 自有 TPU v4 集群 \\
LLaMA 3 405B (2024) & $\sim$54 天 & \$60M--100M（估算） &
  16,384 块 H100，Meta 自有集群 \\
DeepSeek-V3 (2024) & $\sim$2 个月 & \$5.5M（GPU 租赁） &
  2,048 块 H800；不含前期研发成本 \\
GPT-5 / Orion (2025) & 数月 & \$300M--500M（传闻） &
  未经官方确认 \\
下一代 frontier (2026) & --- & \$500M--1B+ &
  行业预估；含多轮训练 \\
\bottomrule
\end{tabularx}
\end{table}

DeepSeek-V3 报告的 \$5.5M 成本
是 AI 行业的一个标志性数字——
它证明了通过工程创新，
训练成本可以比同等性能的竞品低一个数量级。
但需要注意，这个数字仅包含最终训练的 GPU 租赁成本，
不包含数据处理、前期实验、后训练和人员成本。
完整的端到端成本可能在 \$20M--30M 左右。

\subsection{成本的构成}

训练一个 frontier LLM 的总成本可以分解为：

\begin{enumerate}[noitemsep]
  \item \textbf{GPU 计算}（50--70\%）：
        这是最大的成本项。
        以 2026 年初的价格，
        NVIDIA H100 的云端租赁价格约为 \$2.00--3.50/GPU$\cdot$小时，
        新一代 B200 约为 \$4.00--6.00/GPU$\cdot$小时。
        一次使用 8,192 块 H100 训练 60 天的成本：
        $8192 \times 60 \times 24 \times 2.50 \approx \$29.5M$。
  \item \textbf{数据获取与处理}（5--15\%）：
        网页爬取基础设施、数据标注
        （高质量 SFT 数据的标注成本约 \$15--50/对话）、
        数据清洗的计算成本。
  \item \textbf{人员}（15--25\%）：
        AI 研究员的年薪在硅谷通常为 \$300K--800K
        （含股权），顶尖人才超过 \$1M。
        一个 30 人的核心训练团队的年人员成本可达 \$15M--20M。
  \item \textbf{基础设施}（5--10\%）：
        网络设备、存储系统、电力和冷却。
        一个大规模 GPU 集群的年度电力成本可达数百万美元。
  \item \textbf{失败的实验}（不确定）：
        这是最难量化但可能占比很高的成本。
        训练到一半发现数据质量问题、
        超参数配置不佳导致 loss 不收敛——
        这些失败的尝试消耗的计算资源有时比最终成功的训练还多。
\end{enumerate}

\begin{corebox}[2026 年的成本趋势]
两个相反的力量正在塑造训练成本的走向：
\textbf{向上的力量}——模型规模和训练数据量的持续增长，
以及推理模型的 RL 训练需要更多的计算；
\textbf{向下的力量}——更高效的架构
（MoE 的稀疏激活将 FLOPs 降低 5--10 倍）、
更好的硬件（B200 的 FP8 性能是 H100 的 2.5 倍，单位算力价格持续下降）、
更聪明的训练方法（如 DeepSeek 的各种工程优化）。
净效果是：训练同等质量模型的成本在下降，
但 frontier 的定义不断被推高，
所以绝对成本仍在上升。
行业中流传一个说法：
``Training costs are dropping 10x per year for the same quality,
but we keep raising the bar 10x per year.''
\end{corebox}

\subsection{云端 vs.\ 自建集群的经济学}
\label{sec:cloud-vs-onprem}

选择使用云服务商的 GPU 还是自建集群，
是每个 AI 团队面临的关键战略决策。

\textbf{云端（Cloud）的优势}：
无需前期资本支出（CapEx），
按使用量付费（OpEx），
弹性伸缩（训练完成后释放资源），
以及云服务商提供的托管服务
（网络配置、故障自动替换、存储管理）。
主要的云 GPU 提供商包括
AWS（P5 instances with H100/H200, Trn2 instances with Trainium2）、
Google Cloud（A3 with H100, Cloud TPU v5p/v6e）、
Microsoft Azure（ND H100/H200 instances）
和 CoreWeave（GPU 密集型专用云）。
2026 年初，spot instance 价格可以将 H100 成本降低至
\$1.20--1.80/GPU$\cdot$小时，
但 spot instances 可能被中断，
对于需要持续运行数周的预训练任务不太适用。

\textbf{自建集群（On-Premise）的优势}：
长期来看成本更低——
Deloitte 的研究表明，大规模持续使用场景下
自建可以实现云端 60--70\% 的成本。
一块 H100 SXM 的采购价约为 \$25,000--35,000，
如果使用 3 年，每 GPU$\cdot$小时的折旧成本仅约 \$1.00--1.30
（不含电力、冷却、运维人员成本）。
此外，自建集群可以完全控制网络拓扑、
存储架构和安全策略——对于处理高度敏感数据的组织至关重要。

\textbf{关键的交叉点}：
根据 2026 年的 TCO（Total Cost of Ownership）分析，
当 GPU 利用率持续超过 60--70\% 且使用时间超过 18--24 个月时，
自建集群的经济性通常优于云端。
这就是为什么 Meta、Google、Tesla（xAI）等大型玩家
都选择自建——他们的 GPU 几乎 24/7 满载运行。
而初创公司和学术团队通常选择云端——
他们的使用是间歇性的，且无法承担前期的资本投入。

\textbf{混合策略}正在成为主流：
许多中型 AI 公司采用 ``base on-prem + burst to cloud'' 的策略——
自建一个满足日常需求的基础集群，
在需要大规模训练时临时租用云 GPU 扩展容量。

\begin{table}[H]
\centering\small
\caption{云端 vs.\ 自建集群：关键对比（2026 年）}
\begin{tabularx}{\linewidth}{l X X}
\toprule
\textbf{维度} & \textbf{云端} & \textbf{自建} \\
\midrule
前期投入 & 极低 & 数百万--数千万美元 \\
单 GPU 成本/小时 & \$2.00--6.00（按需） & \$1.00--1.80（含运维） \\
弹性 & 分钟级扩缩容 & 数月的采购周期 \\
维护 & 服务商负责 & 需要专业运维团队 \\
数据安全 & 受限于服务商政策 & 完全自主 \\
3 年 TCO（1000 GPU） & \$52M--105M & \$35M--55M \\
适合谁 & 初创、学术、间歇性使用 & 大型实验室、持续训练 \\
\bottomrule
\end{tabularx}
\end{table}

\section{开源训练框架生态}

不同的团队可以选择不同的训练框架，
各有侧重。

\subsection{预训练框架}

\begin{table}[H]
\centering\small
\caption{主流 LLM 预训练框架对比}
\begin{tabularx}{\linewidth}{l l l X}
\toprule
\textbf{框架} & \textbf{维护方} & \textbf{规模} & \textbf{核心优势} \\
\midrule
Megatron-LM & NVIDIA & 产业级 &
  最成熟的 3D 并行实现（TP + PP + DP），
  大多数主要实验室的训练起点 \\
DeepSpeed   & Microsoft & 产业级 &
  ZeRO 优化器~\cite{rajbhandari2020zero}，
  MoE 支持，与 HuggingFace 集成良好 \\
PyTorch FSDP & Meta (PyTorch) & 中大规模 &
  原生 PyTorch API，使用门槛低，
  但在超大规模上不如 Megatron 优化彻底 \\
TorchTitan  & Meta (PyTorch) & 中大规模 &
  LLaMA 训练的参考实现，
  基于 DTensor 分布式原语 \\
Nanotron    & HuggingFace & 研究级 &
  代码简洁可读，适合理解训练流程，
  教学和研究用途 \\
LitGPT      & Lightning AI & 中小规模 &
  易用的训练框架，
  支持多种模型架构，上手快 \\
\bottomrule
\end{tabularx}
\end{table}

\textbf{Megatron-LM} 是大规模预训练的事实标准。
它实现了完整的 3D parallelism（tensor、pipeline、data parallelism），
以及 sequence parallelism 和 expert parallelism 等扩展。
几乎所有训练过 100B+ 模型的团队都在 Megatron 的基础上构建。
它的缺点是代码复杂度高、学习曲线陡峭——
但对于需要在数千块 GPU 上稳定运行数月的训练任务，
没有更可靠的选择。

\textbf{DeepSpeed} 的核心贡献是 ZeRO 优化器，
它通过分片 optimizer states、gradients 和 parameters
大幅降低了每块 GPU 的内存需求。
DeepSpeed-Chat 提供了完整的 RLHF 训练流程支持。
在微调场景下，DeepSpeed 的 ZeRO-Offload
甚至允许在单块 GPU 上微调数十亿参数的模型
（通过将部分状态 offload 到 CPU 内存）。

\textbf{PyTorch FSDP}（Fully Sharded Data Parallel）
是 PyTorch 原生的分布式训练方案，
概念上类似于 DeepSpeed ZeRO Stage 3。
它的优势是与 PyTorch 生态无缝集成，
对于已有 PyTorch 代码库的团队迁移成本最低。

\subsection{后训练/微调框架}

对于后训练（SFT/RLHF/DPO），
生态系统更加多样化：

\begin{itemize}[noitemsep]
  \item \textbf{TRL}（HuggingFace）：
        集成 SFT、DPO、PPO、GRPO 训练，
        与 Transformers 和 PEFT 无缝对接。
        社区最活跃，文档最完善。
  \item \textbf{OpenRLHF}：
        专注于 RLHF/GRPO 训练，
        支持 Ray + vLLM 的大规模分布式 RL。
        在需要将训练规模扩展到数百块 GPU 时，
        OpenRLHF 的分布式 RL 架构更有优势。
  \item \textbf{LLaMA-Factory}：
        一键式微调工具，支持 100+ 模型，
        提供 Web UI 和命令行两种方式。
        特别适合快速实验和原型验证。
  \item \textbf{ms-swift}（阿里巴巴 ModelScope）：
        覆盖 SFT、RL、推理、部署的全流程工具链。
        与中国的模型生态（Qwen、通义千问等）集成最好。
  \item \textbf{Axolotl}：
        微调专用工具，支持多种方法
        （LoRA~\cite{hu2021lora}、QLoRA~\cite{dettmers2023qlora}、
        full fine-tuning、FSDP）。
        配置文件驱动，灵活度高。
  \item \textbf{Unsloth}：
        专注于极致的微调速度和内存优化。
        通过手写 Triton kernel 和自定义 autograd，
        在相同硬件上可以比标准实现快 2--5 倍。
\end{itemize}

\section{开源社区生态：从工具到运动}
\label{sec:opensource-ecosystem}

开源 LLM 不仅是一种技术选择，
更是一场改变 AI 领域权力结构的运动。
2025--2026 年，开源社区的成熟度和影响力达到了新的高度。

\subsection{HuggingFace：ML 领域的 GitHub}

HuggingFace 已经成为 ML 领域的核心基础设施，
类似于 GitHub 之于软件开发。
截至 2026 年初，其平台托管了超过 100 万个模型和
30 万个数据集，月活跃用户超过 500 万。

其生态系统包括：

\begin{itemize}[noitemsep]
  \item \textbf{Model Hub}：托管数十万个预训练模型，
        支持一行代码下载和使用。
        模型卡（Model Card）标准化了模型文档，
        包括训练数据、评估结果、已知限制和许可证信息。
  \item \textbf{Datasets}：标准化的数据集托管和加载，
        覆盖数千个 NLP/CV 数据集。
        FineWeb（15T tokens 的高质量网页数据）
        和 FineWeb-Edu（教育质量过滤的子集）
        是 HuggingFace 对开源训练数据的重要贡献。
  \item \textbf{Transformers 库}：
        统一的 API 支持几乎所有主流模型架构，
        是学术界和工业界最广泛使用的 NLP 库。
        新模型发布后通常在数天内就会被集成。
  \item \textbf{TRL + PEFT}：
        后训练和参数高效微调的标准工具。
        TRL 在 2025 年加入了 GRPO 支持，
        使得复现 DeepSeek-R1 的训练流程成为可能。
  \item \textbf{SmolLM 系列}：
        HuggingFace 自研的小模型家族
        （135M、360M、1.7B 参数），
        展示了数据质量在小模型训练中的关键作用。
        SmolLM2~\cite{allal2025smollm2} 通过精心策划的训练数据
        （Cosmopedia 合成教科书、FineWeb-Edu、Stack 代码数据）
        在同等规模下达到了领先性能。
  \item \textbf{Open LLM Leaderboard}：
        自动化的模型评估排行榜，
        是社区衡量模型质量的重要参考。
        2025 年升级为 v2 版本，
        引入了更难的评估基准（GPQA、MUSR、BBH）
        以应对模型能力的快速提升。
  \item \textbf{Spaces}：
        免费托管模型 demo 的平台，
        研究者可以零成本展示自己的工作。
\end{itemize}

HuggingFace 的核心贡献是\textbf{降低了 ML 的准入门槛}。
在 HuggingFace 之前，使用一个新模型需要：
阅读论文 → 找到作者的 GitHub repo → 配置环境 → 下载权重
→ 编写推理代码。
在 HuggingFace 之后，这简化为三行 Python 代码。

\subsection{本地推理生态}

llama.cpp 和 GGUF 格式的出现
创造了一个完整的本地推理生态系统。
llama.cpp 使用纯 C/C++ 实现，
支持 CPU、Apple Silicon（Metal）、CUDA、Vulkan 等多种后端，
使得在没有 NVIDIA GPU 的消费级硬件上运行 LLM 成为可能。

\textbf{Ollama} 和 \textbf{LM Studio}
进一步将本地推理的门槛降到了非技术用户可及的水平——
一键安装、一键下载模型、开箱即用的聊天界面。
这使得「人人都能在自己的电脑上运行 AI」
从技术理想变成了现实。

社区围绕量化模型建立了活跃的生态：
模型发布后通常在数小时内就会出现
GGUF 格式的各种量化版本
（Q2\_K、Q4\_K\_M、Q5\_K\_S、Q8\_0 等），
用户可以根据自己的硬件条件选择合适的量化级别。

\subsection{社区模型：百花齐放}

开源社区的模型产出已经形成了一个完整的层次结构：

\textbf{第一梯队：实验室级开源}——
DeepSeek-V3/R1（MIT 许可证）、
LLaMA 4~\cite{dubey2024llama3}（Llama 许可证）、
Qwen 3~\cite{qwen2025qwen3}（Apache 2.0）、
Gemma 3（Google，开放权重）。
这些模型的质量已经接近甚至达到闭源模型的水平。

\textbf{第二梯队：完全开源}——
OLMo 2（AI2，数据+代码+权重全部公开）、
SmolLM2（HuggingFace，完整训练流程可复现~\cite{allal2025smollm2}）、
MAP-Neo（M-A-P 社区，7B 全开源 bilingual 模型）。
这些项目的核心价值不在于模型性能的绝对领先，
而在于\textbf{可复现性}——
研究者可以在完全相同的条件下重复实验。

\textbf{第三梯队：社区微调}——
基于第一梯队的 base model，
社区产出了大量针对特定场景的微调模型。
仅在 HuggingFace 上，
基于 LLaMA 的微调模型就超过 10 万个。
这种长尾效应使得开源生态的实际影响力
远超任何单一模型。

\begin{whybox}[开源 LLM 的「飞轮效应」]
开源 LLM 生态形成了一个自我强化的飞轮：
更好的模型 → 更多的用户 → 更多的微调和改进
→ 更活跃的社区 → 更好的工具和基础设施
→ 更低的使用门槛 → 更多的用户 → \ldots
HuggingFace、llama.cpp、Ollama 等工具的成功
证明了这个飞轮的力量。
一个关键观察是：
开源模型的性能与闭源模型的差距
从 2023 年的「代际差距」
缩小到 2025--2026 年的「几个月的差距」——
DeepSeek-R1 在某些数学推理任务上
甚至超越了同期的闭源竞品。
\end{whybox}

\subsection{Chatbot Arena：社区驱动的评估}

LMSYS 的 Chatbot Arena 是开源社区
对 LLM 评估最重要的贡献之一。
它使用类似于国际象棋 Elo 评分的系统，
让真实用户对不同模型的回答进行盲评。
由于评估者不知道哪个回答来自哪个模型，
且样本量足够大（数百万次投票），
Chatbot Arena 的排名被认为是
最能反映真实用户偏好的评估方法——
比任何单一 benchmark 都更可信。

\section{从论文到产品：被低估的工程挑战}

在学术界和媒体报道中，
人们往往关注算法创新和基准测试分数，
而忽略了将 LLM 从实验室推向产品所需的大量工程工作。

\subsection{训练基础设施的可靠性}

在 2048 块 GPU 上训练两个月，
意味着 $2048 \times 60 \times 24 = 295$ 万 GPU$\cdot$小时。
如果每块 GPU 每天有 0.1\% 的故障率，
那么平均每小时会有 $\sim$2 块 GPU 出现故障。
训练系统必须能够自动检测故障、
隔离故障节点、从最近的检查点恢复训练。
DeepSeek-V3~\cite{deepseekv3} 的训练报告强调他们实现了
``zero irrecoverable loss spikes''——
这不是运气，而是精心设计的容错机制的结果。

\subsection{调试分布式训练}

当 loss 在第 50,000 步突然发散，而你用了 4096 块 GPU，
怎么找到 bug？
这是分布式训练中最令人恐惧的场景之一。
可能的原因：某块 GPU 的 ECC 内存错误导致梯度异常；
某个数据 shard 包含异常样本（如超长文本或乱码）；
通信超时导致部分梯度丢失；
甚至可能是浮点数非确定性——
相同的 checkpoint 在不同的 GPU 配置下恢复训练，
可能产生不同的 loss 曲线。

实践中的调试策略包括：
详尽的日志系统（记录每一步的 loss、gradient norm、
learning rate、每块 GPU 的计算时间）；
``canary'' 训练（在小集群上用相同配方并行运行，
比较 loss 曲线是否一致）；
以及最重要的——\textbf{可重现性}，
或者更准确地说，对不可重现性的接受和管理。

\subsection{数据管道的效率}

2048 块 GPU 的训练每秒消耗数百万 token，
数据加载不能成为瓶颈。
这需要高效的数据格式（如 memory-mapped 二进制文件）、
多级缓存（SSD → CPU 内存 → GPU 内存）
和异步预取（在 GPU 计算时提前准备下一批数据）。

\subsection{评估的自动化}

每次 checkpoint 需要在数十个基准测试上评估，
每个测试可能需要数百到数千次推理调用。
这本身就需要大量 GPU 资源。
高效的评估管道（如 lm-evaluation-harness）
和评估集群的合理调度至关重要。

\subsection{人的因素}

最后一个常被忽略的挑战是\textbf{人}。
一次 frontier model 的训练可能消耗数千万美元的计算资源。
如果训练在第六周因为一个 bug 失败而需要重启，
这不仅是经济损失，也是巨大的心理压力。
负责训练的工程师需要 7$\times$24 小时 on-call——
凌晨 3 点收到 loss spike 的告警是常态。
燃尽（burnout）是 LLM 训练团队面临的真实挑战。

这些工程挑战解释了为什么
只有少数组织能够成功训练 frontier LLM——
即使算法和数据都是公开的，
将它们在大规模上可靠地运行仍然是一项极其困难的工程任务。

\section{2025--2026 年的人才与格局变迁}

2025 年下半年至 2026 年初，
AI 领域的人才流动和组织重组以前所未有的烈度展开。
两个标志性事件勾勒出了中国 AI 竞争的新轴线。

\subsection{姚顺雨加入腾讯}

2025 年 12 月 17 日，27--28 岁的\textbf{姚顺雨}（Shunyu Yao）
被任命为腾讯首席 AI 科学家（Chief AI Scientist），
直接向总裁刘炽平汇报。
姚顺雨此前是 OpenAI 研究员，
以 Tree of Thoughts、ReAct 等工作闻名，
是 Agent 领域最具影响力的年轻研究者之一。

腾讯为此次任命做了大规模的组织重构，
同步新设三个一级部门：
AI 基础设施部、AI 数据部和数据算力平台部。
姚顺雨直接统管 AI 基础设施部和大模型部，
获得了罕见的跨部门资源调度权。
在此之前，腾讯已于 12 月 8 日发布了
\textbf{混元 2.0}（Hunyuan 2.0）——
一个 406B/32B MoE 模型，支持 256K 上下文，
为姚顺雨的到来做了技术铺垫。

姚顺雨的研究哲学值得关注：
他多次强调「\textbf{方法的复杂度应匹配任务的难度}」
（method complexity should match task difficulty）——
这是一种反对 benchmark 军备竞赛、
追求方法优雅性的研究品味。
他规划的下一代模型将以 Agent 能力为核心，
参数规模约 30B——
刻意选择了效率优先而非规模优先的路径。
朗桥证券的分析报告指出：
「姚顺雨的蜜月期可能只有 6 个月」——
腾讯需要在 2026 年下半年看到明确的技术突破，
否则这场高调引进的战略价值将面临质疑。

\subsection{梁文锋与 DeepSeek 团队}

如果说 2025 年 1 月的 DeepSeek-R1 是「中国的 Sputnik 时刻」，
那么 2025 年 9 月 R1 登上 \textit{Nature} 封面
则是对这一成就的正式学术认证。
\textbf{梁文锋}（Wenfeng Liang）作为通讯作者，
以不到 200 人的团队完成了 frontier 级别的研究突破。

更令人惊叹的是 DeepSeek 团队在 6 个月内的论文产出：
R1（\textit{Nature}，2025 年 9 月）、
mHC（arXiv:2512.24880，2025 年 12 月 31 日，梁文锋共同作者）、
Engram（arXiv，2026 年 1 月 12 日，梁文锋共同作者）——
三篇论文覆盖了推理训练、架构创新和记忆系统三个方向，
展示了小团队在研究广度上的惊人能力。
DeepSeek APP 的月活用户已突破 1 亿。

V4 的开发据报道聚焦于三个方向：
原生多模态、LTM（Long-Term Memory）突破、
以及国产芯片（华为 Ascend）适配优化。
但梁文锋面临的压力也是空前的——
在 R1 引发全球关注之后，
DeepSeek 必须再次交出
「世界最好的开源模型」这一答卷，
而这一次竞争对手们已经做好了充分准备。

\subsection{竞争格局的两极化}

梁文锋与姚顺雨成为 2026 年中国 AI 的双焦点，
但两人面临的挑战截然不同。

梁文锋是「所有人都在等待他交出下一个作品的明星创业者」——
DeepSeek 必须证明 R1 的成功不是偶然，
V4 需要在技术上超越已经高度内卷的竞争对手，
同时应对来自出口管制和芯片限制的地缘政治压力。

姚顺雨是「从硅谷空降、被委以改革重任的 95 后明星科学家」——
他需要在腾讯这个以产品见长但在 AI 基础研究上相对滞后的组织中，
建立全新的技术文化，
在政治资本耗尽之前拿出让人信服的成果。

两种模式——
小团队极致效率 vs.\ 巨头组织重构——
谁能在 2026 年下半年率先交出答卷，
将深刻影响中国乃至全球 AI 产业的走向。

\subsection{全球 AI 产品生态的演变}

在模型层之外，AI 产品层的竞争同样激烈。

\textbf{Cognition AI} 在 2025 年底收购了 Windsurf IDE，
将 Devin 的自主编程能力与 IDE 的交互体验融合，
试图定义「AI 原生开发环境」的新形态。
这一收购代表了 Agent 领域的一个重要趋势：
从纯自主模式向人机协作模式的收敛——
最有效的 AI 编程工具可能既不是纯自主 Agent（Devin），
也不是纯辅助工具（Copilot），
而是两者的有机结合。


%% ============================================================
%% BibTeX entries referenced in this chapter
%% (to be added to the main .bib file)
%%
%% @article{deepseekv3,
%%   title   = {{DeepSeek-V3} Technical Report},
%%   author  = {DeepSeek-AI},
%%   journal = {arXiv preprint arXiv:2412.19437},
%%   year    = {2024},
%% }
%%
%% @article{deepseekr1,
%%   title   = {{DeepSeek-R1}: Incentivizing Reasoning Capability in {LLMs} via Reinforcement Learning},
%%   author  = {DeepSeek-AI},
%%   journal = {arXiv preprint arXiv:2501.12948},
%%   year    = {2025},
%% }
%%
%% @article{penedo2024fineweb,
%%   title   = {{FineWeb}: decanting the web for the finest text data at scale},
%%   author  = {Penedo, Guilherme and Kydl{\'\i}{\v{c}}ek, Hynek and {Ben Allal}, Loubna and others},
%%   journal = {arXiv preprint arXiv:2406.17557},
%%   year    = {2024},
%% }
%%
%% @article{kaplan2020scaling,
%%   title   = {Scaling Laws for Neural Language Models},
%%   author  = {Kaplan, Jared and McCandlish, Sam and Henighan, Tom and others},
%%   journal = {arXiv preprint arXiv:2001.08361},
%%   year    = {2020},
%% }
%%
%% @article{hoffmann2022chinchilla,
%%   title   = {Training Compute-Optimal Large Language Models},
%%   author  = {Hoffmann, Jordan and Borgeaud, Sebastian and Mensch, Arthur and others},
%%   journal = {arXiv preprint arXiv:2203.15556},
%%   year    = {2022},
%% }
%%
%% @article{dubey2024llama3,
%%   title   = {The {Llama 3} Herd of Models},
%%   author  = {Dubey, Abhimanyu and Jauhri, Abhinav and Pandey, Abhinav and others},
%%   journal = {arXiv preprint arXiv:2407.21783},
%%   year    = {2024},
%% }
%%
%% @article{bai2022constitutional,
%%   title   = {Constitutional {AI}: Harmlessness from {AI} Feedback},
%%   author  = {Bai, Yuntao and Kadavath, Saurav and Kundu, Sandipan and others},
%%   journal = {arXiv preprint arXiv:2212.08073},
%%   year    = {2022},
%% }
%%
%% @article{touvron2023llama,
%%   title   = {{LLaMA}: Open and Efficient Foundation Language Models},
%%   author  = {Touvron, Hugo and Lavril, Thibaut and Izacard, Gautier and others},
%%   journal = {arXiv preprint arXiv:2302.13971},
%%   year    = {2023},
%% }
%%
%% @article{rajbhandari2020zero,
%%   title   = {{ZeRO}: Memory Optimizations Toward Training Trillion Parameter Models},
%%   author  = {Rajbhandari, Samyam and Rasley, Jeff and Ruwase, Olatunji and He, Yuxiong},
%%   journal = {arXiv preprint arXiv:1910.02054},
%%   year    = {2019},
%% }
%%
%% @article{hu2021lora,
%%   title   = {{LoRA}: Low-Rank Adaptation of Large Language Models},
%%   author  = {Hu, Edward J and Shen, Yelong and Wallis, Phillip and others},
%%   journal = {arXiv preprint arXiv:2106.09685},
%%   year    = {2021},
%% }
%%
%% @article{dettmers2023qlora,
%%   title   = {{QLoRA}: Efficient Finetuning of Quantized Language Models},
%%   author  = {Dettmers, Tim and Pagnoni, Artidoro and Holtzman, Ari and Zettlemoyer, Luke},
%%   journal = {arXiv preprint arXiv:2305.14314},
%%   year    = {2023},
%% }
%%
%% @article{jiang2024mixtral,
%%   title   = {Mixtral of Experts},
%%   author  = {Jiang, Albert Q and Sablayrolles, Alexandre and Roux, Antoine and others},
%%   journal = {arXiv preprint arXiv:2401.04088},
%%   year    = {2024},
%% }
%%
%% @article{anil2023gemini,
%%   title   = {Gemini: A Family of Highly Capable Multimodal Models},
%%   author  = {Anil, Rohan and Dai, Andrew M and Firat, Orhan and others},
%%   journal = {arXiv preprint arXiv:2312.11805},
%%   year    = {2023},
%% }
%%
%% @article{qwen2025qwen3,
%%   title   = {{Qwen3} Technical Report},
%%   author  = {Qwen Team},
%%   journal = {arXiv preprint},
%%   year    = {2025},
%% }
%%
%% @article{allal2025smollm2,
%%   title   = {{SmolLM2}: When Smol Goes Big --- Data-Centric Training of a Small Language Model},
%%   author  = {{Ben Allal}, Loubna and Lozhkov, Anton and Bakouch, Elie and others},
%%   journal = {arXiv preprint arXiv:2502.02737},
%%   year    = {2025},
%% }
